\section{Trả lại thứ tự cho câu}
\label{sec:ch07-vi-tri}

\index{mã hóa vị trí!bài toán}%
Mục~\ref{sec:ch07-tu-nhin} đo được rằng tầng attention không phân biệt nổi hai
thứ tự, sai khác cỡ \(4 \times 10^{-8}\). Mục này là chỗ vá.

Cách vá của bài rẻ tới mức đáng ngờ: cộng thêm vào embedding một vector phụ
thuộc vị trí. Không phải nối vào, mà cộng vào, nên nó không tốn thêm chiều nào.
Bài gọi là \tn{mã hóa vị trí}{positional encoding} và chọn dạng sin cos:

\begin{equation}
  PE_{(pos, 2i)} = \sin\!\left( \frac{pos}{10000^{2i/\dmodel}} \right),
  \qquad
  PE_{(pos, 2i+1)} = \cos\!\left( \frac{pos}{10000^{2i/\dmodel}} \right).
  \label{eq:ch07-pe}
\end{equation}

Chỉ số \(i\) chạy trên các \emph{cặp} cột, không phải trên từng cột. Cột \(2i\)
và cột \(2i+1\) dùng chung một tần số, sin ở cột chẵn và cos ở cột lẻ. Nên bảng
này không phải \(\dmodel\) đường cong độc lập, nó là \(\dmodel/2\) cặp quay, và
toàn bộ mục này chạy trên chuyện đó.

\begin{minted}{python}
position = torch.arange(n_positions, dtype=torch.float32).unsqueeze(1)
pair = torch.arange(0, d_model, 2, dtype=torch.float32)
angle = position / torch.pow(PE_BASE, pair / d_model)
table = torch.zeros(n_positions, d_model)
table[:, 0::2] = torch.sin(angle)
table[:, 1::2] = torch.cos(angle)
\end{minted}

\subsection*{Lý do bài đưa ra, và chữ bài dùng}

\index{mã hóa vị trí!bài nói phỏng đoán}%
Mục 3.5 giải thích lựa chọn bằng một câu:

\begin{quote}
\enquote{We chose this function because we hypothesized it would allow the
model to easily learn to attend by relative positions, since for any fixed
offset \(k\), \(PE_{pos+k}\) can be represented as a linear function of
\(PE_{pos}\).}
\end{quote}

Chữ \enquote{hypothesized} là chữ của bài, và nó đáng để ý theo cùng cách chữ
\enquote{conjecture} của Bahdanau đáng để ý ở
mục~\ref{sec:ch06-ai-do}. Bài không chứng minh tính chất ấy, không đo nó, và
không có thí nghiệm nào trong bài tách được nó ra khỏi phần còn lại.

Nên tôi đo. Phần toán thì ngắn: với một cặp cột ở tần số
\(\omega_i = 1/10000^{2i/\dmodel}\), dịch vị trí đi \(k\) là quay cặp ấy một góc
\(k\omega_i\), vì

\begin{equation}
  \begin{bmatrix} \sin((pos+k)\omega_i) \\ \cos((pos+k)\omega_i) \end{bmatrix}
  =
  \begin{bmatrix} \cos(k\omega_i) & \sin(k\omega_i) \\
                 -\sin(k\omega_i) & \cos(k\omega_i) \end{bmatrix}
  \begin{bmatrix} \sin(pos\,\omega_i) \\ \cos(pos\,\omega_i) \end{bmatrix},
  \label{eq:ch07-quay}
\end{equation}

đúng bằng công thức cộng góc. Ghép \(\dmodel/2\) khối ấy lại theo đường chéo thì
được \(\mat{M}_k\). Nội dung thật của phát biểu nằm ở chỗ \(pos\) không xuất
hiện ở vế phải của ma trận: cùng một \(\mat{M}_k\) dùng được ở mọi vị trí.

\begin{measured}{\(\mat{M}_k \, PE_{pos}\) so với \(PE_{pos+k}\), \(\dmodel = 64\)}
\begin{minted}{text}
offset k  max|M_k PE(pos) - PE(pos+k)|  over positions
1         1.341e-06                    5 tested
2         1.609e-06                    5 tested
5         1.222e-06                    5 tested
13        1.840e-06                    5 tested
40        2.623e-06                    5 tested
\end{minted}

\(\mat{M}_k\) dựng bằng giải tích theo \eqref{eq:ch07-quay}, không khớp số.
Các vị trí thử là 0, 1, 7, 30, 54.

Đo bằng \code{experiments/ch07\_position.py} tại tag \code{ch07}.
\end{measured}

Sai khác cỡ \(10^{-6}\) trên mọi cặp \((k, pos)\). Phỏng đoán của bài đúng.

\index{mã hóa vị trí!khớp M\_k bằng bình phương tối thiểu thì hỏng}%
Một chỗ bẫy, ghi ra vì nó chờ sẵn đúng người đọc chịu tự kiểm. Cách tự nhiên
nhất để kiểm \enquote{tồn tại một ánh xạ tuyến tính} là đi \emph{khớp} nó: lấy
bình phương tối thiểu trên một cửa sổ vị trí rồi xem \(\mat{M}_k\) thu được có
đổi theo \(pos\) không. Làm thế thì kết quả trông như thể nó có đổi. Không phải
vậy. Các chiều bước sóng dài gần như không quay trên một trăm vị trí, nên ma
trận thiết kế suy biến nặng và bài toán khớp không xác định được ở những chiều
ấy. Dựng \(\mat{M}_k\) bằng giải tích thì không có vấn đề đó, và bảng trên là
cách dựng ấy.

\subsection*{Một tính chất mạnh hơn cái bài phát biểu}

\index{mã hóa vị trí!tích vô hướng chỉ phụ thuộc khoảng cách}%
Trong lúc đo, tôi thử một đại lượng khác và nó cho một phát biểu sạch hơn.

Lấy \tn{tích vô hướng}{dot product} của hai hàng cách nhau \(k\). Mỗi cặp cột
đóng góp
\(\sin(a)\sin(b) + \cos(a)\cos(b) = \cos(a - b)\), và hiệu \(a - b\) đúng bằng
\(k\omega_i\) bất kể \(pos\). Cộng lại:

\begin{equation}
  PE_{pos} \cdot PE_{pos+k} = \sum_{i=0}^{\dmodel/2 - 1} \cos(k \omega_i).
  \label{eq:ch07-tich-vo-huong-pe}
\end{equation}

Vế phải không có \(pos\).

\begin{measured}{\(PE_{pos} \cdot PE_{pos+k}\) tại các vị trí khác nhau, \(\dmodel = 64\)}
\begin{minted}{text}
offset k  pos=0     pos=1     pos=7     pos=30    spread
1         30.91683  30.91683  30.91683  30.91683  0.00e+00
2         28.30386  28.30386  28.30386  28.30386  1.91e-06
5         23.50397  23.50397  23.50397  23.50397  0.00e+00
13        21.00672  21.00672  21.00672  21.00672  3.81e-06
40        15.39746  15.39746  15.39746  15.39746  3.81e-06
\end{minted}

Đo bằng \code{experiments/ch07\_position.py} tại tag \code{ch07}.
\end{measured}

Cột \code{spread} bằng 0 hoặc bằng nhiễu float. Bài chỉ nói tồn tại một ánh xạ
tuyến tính; đây nói tích vô hướng của hai vị trí là một hàm của riêng khoảng
cách giữa chúng. Vì attention chấm điểm \emph{bằng} tích vô hướng, đó đúng là
đại lượng cơ chế nhìn thấy.

\subsection*{Nhưng đừng đi xa hơn thế, và đây là chỗ cách kể thông thường sai}

\index{mã hóa vị trí!tích vô hướng không đơn điệu theo k}%
Từ bảng trên, bước tiếp theo nghe rất xuôi tai: tích giảm khi khoảng cách tăng,
30.91683 rồi 28.30386 rồi 23.50397 rồi 21.00672 rồi 15.39746, nên mô hình đọc
ngược ra được khoảng cách. Tôi đã định viết đúng câu ấy.

Nó sai, và sai ở chỗ đo được.

\begin{measured}{chỗ \(PE_{pos} \cdot PE_{pos+k}\) quay đầu}
\begin{minted}{text}
d_model  k where it first rises again  value at k=1  min over k<200
32       5                             15.3136       2.2507 at k=161
64       6                             30.9168       7.0883 at k=172
512      44                            249.1021      86.3717 at k=199
\end{minted}

Đo bằng \code{experiments/ch07\_position.py} tại tag \code{ch07}.
\end{measured}

Tổng các cosin ở \eqref{eq:ch07-tich-vo-huong-pe} không đơn điệu. Ở
\(\dmodel = 64\) nó giảm tới \(k = 5\) rồi tăng lại ở \(k = 6\); ở
\(\dmodel = 32\) chỗ quay đầu là \(k = 5\); ở \(\dmodel = 512\) là \(k = 44\).
Sau chỗ ấy, một giá trị của tích ứng với nhiều khoảng cách, nên không có phép
đọc ngược nào cả.

Chỗ này đáng đọc cẩn thận, vì tôi \emph{không} đang bắt lỗi bài. Bài phát biểu
đúng một chuyện, là tính chất tuyến tính, và chuyện ấy đúng. Cái sai nằm ở bước
suy diễn thêm mà cách kể lại hay gắn vào. Cùng loại với chuyện
chương~\ref{ch:encoder-decoder} và chương~\ref{ch:mot-vector-cho-moi-tu} phải gỡ
về \enquote{encoder-decoder hỏng ở câu dài}: bài không nói thế, người kể lại nói
thế.

Và nó không phải chuyện xa xôi. Sách này chạy \(\dmodel\) 24 tới 32 với câu dài
nhất 23 token, tức nằm hẳn bên phải chỗ quay đầu. Ở quy mô của bài,
\(\dmodel = 512\) và câu 15 tới 40 từ, phần lớn khoảng cách nằm bên trái
\(k = 44\), nên ở đó cách đọc kia gần đúng. Một tính chất phụ thuộc quy mô mà
người ta trích như một tính chất tuyệt đối.

\subsection*{Dải bước sóng, và một câu của bài không đúng theo nghĩa đen}

\index{mã hóa vị trí!dải bước sóng thật}%
Bài mô tả tập tần số bằng một câu: \enquote{The wavelengths form a geometric
progression from \(2\pi\) to \(10000 \cdot 2\pi\)}.

Cấp số nhân thì đúng. Hai đầu mút thì không hẳn. Bước sóng của cặp \(i\) là
\(2\pi \cdot 10000^{2i/\dmodel}\), và \(2i\) lớn nhất là \(\dmodel - 2\), nên
bước sóng dài nhất công thức đạt tới là \(2\pi \cdot 10000^{(\dmodel-2)/\dmodel}\),
luôn ngắn hơn \(10000 \cdot 2\pi\).

\begin{measured}{bước sóng dài nhất thật sự đạt tới, tính theo đơn vị \(2\pi\)}
\begin{minted}{text}
d_model  shortest/2pi  longest/2pi  paper says  shortfall
64       1.0000        7498.94      10000.0     1.3335x
128      1.0000        8659.64      10000.0     1.1548x
512      1.0000        9646.62      10000.0     1.0366x
\end{minted}

Đo bằng \code{experiments/ch07\_position.py} tại tag \code{ch07}.
\end{measured}

Ở \(\dmodel = 512\), tức cấu hình của bài, chênh 1.0366 lần và câu của bài coi
như đúng. Ở \(\dmodel = 64\) chênh 1.3335 lần. Câu ấy đúng theo nghĩa tiệm cận
khi \(\dmodel\) lớn, không đúng theo nghĩa đen, và khoảng cách giữa hai nghĩa
lớn dần khi mô hình nhỏ đi.

\subsection*{Bài tự thử cách khác, và cách khác cũng tốt bằng}

\index{mã hóa vị trí!hàng (E) của bảng 3}%
Chỗ trung thực nhất của mục 3.5 là câu cuối:

\begin{quote}
\enquote{We also experimented with using learned positional embeddings instead,
and found that the two versions produced nearly identical results (see Table 3
row (E)). We chose the sinusoidal version because it may allow the model to
extrapolate to sequence lengths longer than the ones encountered during
training.}
\end{quote}

Hàng (E) cho BLEU 25.7 so với 25.8 của base. Nghĩa là toàn bộ cấu trúc đẹp đẽ ở
trên, cấp số nhân bước sóng, phép quay, tính chất khoảng cách, đổi lấy đúng 0.1
BLEU so với một bảng tham số học được từ đầu.

Lý do bài giữ sin cos không phải kết quả đo được. Nó là một phỏng đoán khác,
rằng dạng này ngoại suy được ra câu dài hơn lúc huấn luyện, và câu ấy dùng chữ
\enquote{may}. Chương~\ref{ch:vi-tri-va-ngu-canh} là chương đo phỏng đoán ấy, và
tôi nói trước kết quả vì nó đổi cách đọc mục này: sinusoid không ngoại suy tốt,
và cái thay thế nó bốn năm sau lấy đúng ý tưởng phép quay ở
\eqref{eq:ch07-quay} rồi đặt phép quay ấy vào chỗ khác.
